\section{Relation to prior institutions and systems}\label{sec:prior}

\paragraph{The guild.} The craft guilds of preindustrial Europe are the nearest historical model for an institution
that holds and transmits working knowledge, and historians disagree about them in an instructive way. On one
account guilds existed to provide transferable skills through apprenticeship, and they spread technique through
migrant labour \citep{epstein1998craft}. On another they excluded competitors, manipulated markets, blocked
innovations and denied apprenticeships to outsiders, to the benefit of their members and of political elites
\citep{ogilvie2019guilds}. Both accounts describe one institution, and we read them together as follows: a guild
could pay for the keeping of knowledge only out of the returns of exclusion. A Bonfire keeps the first function and
has no need of the second. It is paid from issuance for what other communities take up, so it gains by publishing
and loses nothing when anyone uses what it publishes. Membership is open as well. What a guild's master held by
admission, a Bonfire's stakeholder holds by work: stake is earned through contribution.

\paragraph{Issuance that follows opinion, stake or volume.} Bittensor allocates by the scores
validators assign to one another's work, combined by stake. An empirical study of 64 of its sub-networks found that
validators' rewards tracked their stake, with correlations between $0.80$ and $0.95$, while miners' rewards tracked
measured performance only weakly, between $0.10$ and $0.30$ \citep{lui2025bittensor}, and the project's own
documentation states that a validator who copies the expected consensus can earn more per unit of stake than one who
evaluates honestly \citep{bittensor2026weightcopy}. The closest existing design is The Graph. Developers publish
data definitions, staked indexers serve them, and an indexer's signed attestation of a query response can be
disputed and re-run by an arbitrator \citep{thegraph2026docs}. Its issuance, however, follows stake and not use:
indexing rewards are distributed in proportion to the signal curators place on a definition and the stake indexers
allocate to it. The project itself reports that neither quantity measures whether an indexer served a query, that
15.2\% of indexing rewards in 2025 went to indexers that were not delivering value, and that it has added an
eligibility test based on queries served \citep{thegraph2026eligibility}. Where staked votes alone direct emissions,
as with Curve's gauge weights \citep{curve2026gauges}, markets have formed in which those votes are bought
\citep{votium2026docs}. And a marketplace that rewarded reported trading volume saw the volume manufactured: almost
80\% of the wash-trading operations identified on LooksRare closed at a profit \citep{lamorgia2023nft}. These are
the reasons the second agreement is reached by re-execution, where there is nothing to copy; that stake is a
ceiling and not a weight; and that what fills the ceiling is breadth and not volume.

\paragraph{Paying on speech and attention.} Several systems have attached payment to posts or to persons.
Friend.tech let users buy ``keys'' in individuals; it collected almost \$90 million in fees, and its daily fees,
which once reached \$2 million, stood below \$100 by September 2024 \citep{craig2024friendtech}. On Steemit, where
stake-weighted votes on posts directed rewards, more than 16\% of transfers went to accounts suspected of being bots
that sold votes \citep{li2019steem}. On Farcaster, token rewards raised the volume of posting while often failing to
raise, and sometimes lowering, its quality, and 92.4\% of the tipping transactions of one token came from two bot
accounts \citep{yang2025farcaster}. In all three the quantity paid on was a proxy at the level of the person or the
post, and in the last two the proxy was produced by bots. We draw a
narrow conclusion: a use that a third party can re-execute is a different kind of evidence from a vote, a price or a
tip, and it still needs a defence against one party posing as many. For the same reason the protocol pays nothing for
passing a construction on. A reward for transmission is a reward for attention, and it is paid before anyone knows
whether what spread was kept.

\paragraph{One ledger, many languages.} A protocol built on naming games might be expected to drive its
participants toward a single shared language. The literature suggests otherwise. On networks with strong community
structure, naming games settle into clusters that keep different conventions for long periods, in some cases
indefinitely \citep{lu2009naming,dallasta2006nonequilibrium}. A convention can also be captured by a small committed
group: on a complete graph, a committed minority of about 10\% changes the time to consensus from exponential in the
population to logarithmic \citep{xie2011social}. Populations of language-model agents form conventions in
the same way \citep{ashery2025emergent}. We take this to mean that agreement on meaning is neither stable enough nor
secure enough to carry payment. The protocol asks for agreement only on the ledger, and meaning stays plural: breadth counts the communities that
hold a construction and does not require that all of them do.

\paragraph{A studied system of discrete rules.} Collectible card games are a studied case of a population using
discrete rules with definite effects: the cards of one such game are rules precise enough that the game has been
shown to be Turing complete \citep{churchill2021magic}, and what a card contributes depends on what it is combined
with \citep{kritz2025synergy}. The cards trade in a real market \citep{luckingreiley1999magic}, in which cards
designated rare sell for more than functionally equivalent substitutes \citep{hughes2022rarity}. That market
prices scarcity over and above function. A protocol built on open use has no scarcity of that kind to price. What
is scarce in it is credit for having been first (Section~\ref{sec:lineage}).

\paragraph{The scientific process.} Science rewards the structuring of knowledge so that its use produces more of
it, and it is the working precedent for open use with exclusive credit. A result is free to use, and what its
author holds is priority and citation, the counterparts here of registration and lineage. 
